% !TEX program = xelatex
\documentclass[UTF8,zihao=-4,a4paper,fontset=none]{ctexart}

% ================================================================
% 页面设置：A4；Word 默认页边距，上下2.54cm，左右3.17cm
% ================================================================
\usepackage[a4paper,top=2.54cm,bottom=2.54cm,left=3.17cm,right=3.17cm]{geometry}

% ================================================================
% 字体设置：调用 Overleaf 同级目录下上传的字体文件
% ================================================================
\def\FontPath{./}

\usepackage{fontspec}
\usepackage{xeCJK}
\defaultfontfeatures{Ligatures=TeX}

\setmainfont{TeX Gyre Termes}

% 正文宋体：方正书宋简体
\setCJKmainfont[
    Path=\FontPath,
    AutoFakeBold=2
]{方正书宋简体.TTF}

% 中文字体族
\setCJKfamilyfont{song}[
    Path=\FontPath,
    AutoFakeBold=2
]{方正书宋简体.TTF}

\setCJKfamilyfont{hei}[
    Path=\FontPath,
    AutoFakeBold=2
]{方正黑体简体.TTF}

\setCJKfamilyfont{kai}[
    Path=\FontPath,
    AutoFakeBold=2
]{方正楷体简体.TTF}

\setCJKfamilyfont{fs}[
    Path=\FontPath,
    AutoFakeBold=2
]{仿宋_GB2312.ttf}

% 方正小标宋简体：用于封面大标题、论文题目
\newCJKfontfamily\xiaobiaosong[
    Path=\FontPath,
    AutoFakeBold=2
]{方正小标宋简体.ttf}

\newcommand{\FZsong}{\CJKfamily{song}}
\newcommand{\FZhei}{\CJKfamily{hei}}
\newcommand{\FZkai}{\CJKfamily{kai}}
\newcommand{\FZfang}{\CJKfamily{fs}}

% ================================================================
% 正文行距与段落：宋体小四，固定值24磅，段前段后0行，首行缩进2字符
% ================================================================
\usepackage{setspace}
\setlength{\parindent}{2em}
\setlength{\parskip}{0pt}

\newcommand{\maintextformat}{%
    \FZsong\zihao{-4}\setlength{\baselineskip}{24pt}\selectfont
}

% ================================================================
% 标题层级：一、；（一）；1.；①
% ================================================================
\newcommand{\circlednum}[1]{%
    \ifcase#1
    \or ①
    \or ②
    \or ③
    \or ④
    \or ⑤
    \or ⑥
    \or ⑦
    \or ⑧
    \or ⑨
    \or ⑩
    \else #1
    \fi
}

\renewcommand{\theparagraph}{\circlednum{\value{paragraph}}}

\ctexset{
    section = {
        name = {},
        number = \chinese{section}、,
        format = \FZhei\zihao{-3}\setlength{\baselineskip}{24pt}\selectfont,
        beforeskip = 0pt,
        afterskip = 0pt,
        aftername = {}
    },
    subsection = {
        name = {（,）},
        number = \chinese{subsection},
        format = \FZkai\zihao{4}\setlength{\baselineskip}{24pt}\selectfont,
        beforeskip = 0pt,
        afterskip = 0pt,
        aftername = {}
    },
    subsubsection = {
        name = {},
        number = \arabic{subsubsection}.,
        format = \FZsong\bfseries\zihao{-4}\setlength{\baselineskip}{24pt}\selectfont,
        beforeskip = 0pt,
        afterskip = 0pt,
        aftername = {}
    },
    paragraph = {
        name = {},
        number = \theparagraph,
        format = \FZsong\zihao{-4}\setlength{\baselineskip}{24pt}\selectfont,
        beforeskip = 0pt,
        afterskip = 0pt,
        aftername = {}
    }
}

% ================================================================
% 数学公式：可编辑、居中，公式编号右对齐，格式为（1）
% ================================================================
\usepackage{amsmath,amssymb}
\usepackage{unicode-math}
\setmathfont{Latin Modern Math}
\newcommand{\bm}[1]{\symbfit{#1}}
\renewcommand{\theequation}{\arabic{equation}}

% ================================================================
% 图表与浮动体
% ================================================================
\usepackage{graphicx}
\usepackage{tikz}
\usetikzlibrary{arrows.meta,calc,positioning,fit}
\usepackage{booktabs,threeparttable,multirow,float,array}
\usepackage[section]{placeins}
\usepackage[table]{xcolor}
\definecolor{tablehead}{HTML}{F2F5F8}
\usepackage{subcaption}
\usepackage{caption}
\usepackage{etoolbox}

\DeclareCaptionLabelFormat{cnlabel}{#1#2}
\DeclareCaptionLabelSeparator{twospace}{\quad}
\DeclareCaptionFont{songxiaosi}{%
    \FZsong\zihao{-4}\setlength{\baselineskip}{24pt}\selectfont
}

\captionsetup{
    labelformat=cnlabel,
    labelsep=twospace,
    justification=centering,
    singlelinecheck=false,
    font=songxiaosi
}

\captionsetup[table]{name=表,position=top}
\captionsetup[figure]{name=图,position=bottom}

\renewcommand{\thetable}{\arabic{table}}
\renewcommand{\thefigure}{\arabic{figure}}

\renewcommand{\arraystretch}{1.0}
\setlength{\tabcolsep}{6pt}

\AtBeginEnvironment{tabular}{%
    \FZsong\zihao{5}\setlength{\baselineskip}{12pt}\selectfont
}

\setlength{\floatsep}{8pt plus 2pt minus 2pt}
\setlength{\textfloatsep}{10pt plus 2pt minus 3pt}
\setlength{\intextsep}{8pt plus 2pt minus 2pt}
\setlength{\abovecaptionskip}{5pt}
\setlength{\belowcaptionskip}{0pt}

% ================================================================
% 列表
% ================================================================
\usepackage{enumitem}
\setlist{
    nosep,
    leftmargin=2em,
    itemsep=0pt,
    topsep=0pt,
    parsep=0pt,
    partopsep=0pt
}

% ================================================================
% 目录与超链接：目录引用至正文第三级标题
% ================================================================
\usepackage[hidelinks]{hyperref}
\usepackage{tocloft}

\setcounter{secnumdepth}{4}
\setcounter{tocdepth}{3}

\renewcommand{\contentsname}{目\quad 录}
\renewcommand{\cfttoctitlefont}{%
    \hfill\FZhei\zihao{4}\setlength{\baselineskip}{24pt}\selectfont
}
\renewcommand{\cftaftertoctitle}{\hfill}
\setlength{\cftbeforetoctitleskip}{0pt}
\setlength{\cftaftertoctitleskip}{0pt}

\renewcommand{\cftsecfont}{\FZsong\zihao{-4}\setlength{\baselineskip}{24pt}\selectfont}
\renewcommand{\cftsubsecfont}{\FZsong\zihao{-4}\setlength{\baselineskip}{24pt}\selectfont}
\renewcommand{\cftsubsubsecfont}{\FZsong\zihao{-4}\setlength{\baselineskip}{24pt}\selectfont}

\renewcommand{\cftsecpagefont}{\FZsong\zihao{-4}}
\renewcommand{\cftsubsecpagefont}{\FZsong\zihao{-4}}
\renewcommand{\cftsubsubsecpagefont}{\FZsong\zihao{-4}}

\setlength{\cftsecnumwidth}{2.8em}
\setlength{\cftsubsecnumwidth}{3.8em}
\setlength{\cftsubsubsecnumwidth}{3.2em}

\setlength{\cftbeforesecskip}{0pt}
\setlength{\cftbeforesubsecskip}{0pt}
\setlength{\cftbeforesubsubsecskip}{0pt}

% ================================================================
% 表格与插图清单：同一页标题，依次列出表、图
% ================================================================
\makeatletter
\newcommand{\listoftablesandfigures}{%
    \clearpage
    \phantomsection
    \begin{center}
        {\FZhei\zihao{4}\setlength{\baselineskip}{24pt}\selectfont 表格与插图清单\par}
    \end{center}
    \vspace{0pt}
    {\FZsong\zihao{-4}\setlength{\baselineskip}{24pt}\selectfont
    \begingroup
    \let\oldnumberline\numberline
    \renewcommand{\numberline}[1]{表##1\quad}
    \@starttoc{lot}
    \renewcommand{\numberline}[1]{图##1\quad}
    \@starttoc{lof}
    \endgroup}
}
\makeatother

% ================================================================
% 页眉页脚：正文页码从正文第一页开始
% ================================================================
\usepackage{fancyhdr}
\fancypagestyle{plain}{
    \fancyhf{}
    \renewcommand{\headrulewidth}{0pt}
    \renewcommand{\footrulewidth}{0pt}
    \fancyfoot[C]{\FZsong\zihao{-5}\thepage}
}

% ================================================================
% 图片路径
% ================================================================
\graphicspath{{figures/}}

% ================================================================
% 论文基本信息
% 除封面、致谢外，不得出现学校、队员、指导教师信息
% ================================================================
\newcommand{\workID}{TJJM20260000000}
\newcommand{\schoolName}{福州大学}
\newcommand{\paperTitle}{代际友好视角下福州市社区养老服务供需错配识别与智能补位优化研究\\[0.3em]——基于多源 POI、人口栅格与图神经网络}
\newcommand{\paperTitlePlain}{代际友好视角下福州市社区养老服务供需错配识别与智能补位优化研究——基于多源 POI、人口栅格与图神经网络}
\newcommand{\teamMembers}{蔡明润\quad 王智洋\quad 杨孙源}
\newcommand{\advisorName}{XXX}

\newcommand{\keywords}[1]{%
    \par
    \noindent{\FZhei\zihao{-4} 关键词：}{\FZsong\zihao{-4}#1}
    \par
}

% ================================================================
% 封面信息栏：贴着文字的下划线
% ================================================================

% 单行字段：学校、队员、指导老师
\newcommand{\coverfield}[2]{%
    \noindent
    \makebox[3.4cm][r]{%
        {\FZfang\zihao{-2}\setlength{\baselineskip}{24pt}\selectfont #1：}%
    }%
    \hspace{0.45cm}%
    {\FZfang\zihao{-2}\setlength{\baselineskip}{24pt}\selectfont
    \underline{\makebox[9.2cm][c]{#2}}}%
    \par\vspace{0.95cm}%
}

% 多行字段：论文题目
\newcommand{\covertitlefield}{%
    \noindent
    \makebox[3.4cm][r]{%
        {\FZfang\zihao{-2}\setlength{\baselineskip}{24pt}\selectfont 论文题目：}%
    }%
    \hspace{0.45cm}%
    \begin{minipage}[t]{9.2cm}
        \centering
        {\FZfang\zihao{-2}\setlength{\baselineskip}{32pt}\selectfont
        \underline{\makebox[9.2cm][c]{代际友好视角下福州市社区养老服务}}\par
        \underline{\makebox[9.2cm][c]{供需错配识别与智能补位优化研究}}\par
        \underline{\makebox[9.2cm][c]{——基于多源 POI、人口栅格与图神经网络}}\par}
    \end{minipage}
    \par\vspace{0.95cm}%
}

\begin{document}
\maintextformat
\pagestyle{empty}
\pagenumbering{gobble}

% ================================================================
% 封面页：严格按封面示例；仅封面可出现学校、参赛队员、指导教师信息
% ================================================================
\begin{titlepage}
    \maintextformat

    % 作品编号：黑体三号，位于左上
    \noindent{\FZhei\zihao{3}\setlength{\baselineskip}{24pt}\selectfont 作品编号：\workID\par}

    \vspace{2.1cm}

    % 大赛名称与参赛作品：方正小标宋简体一号，固定45磅
    \begin{center}
        {\xiaobiaosong\zihao{1}\setlength{\baselineskip}{45pt}\selectfont
        2026年（第十二届）全国大学生统计建模大赛\par
        参\quad 赛\quad 作\quad 品\par}
    \end{center}

    \vspace{2.7cm}

    % 封面信息：仿宋_GB2312 小二号，文字贴着下划线
    \begin{center}
    \begin{minipage}{13.2cm}
        \coverfield{参赛学校}{\schoolName}
        \covertitlefield
        \coverfield{参赛队员}{\teamMembers}
        \coverfield{指导老师}{\advisorName}
    \end{minipage}
    \end{center}

\end{titlepage}
\clearpage

% ================================================================
% 摘要与关键词
% ================================================================
\begin{center}
    {\xiaobiaosong\zihao{3}\setlength{\baselineskip}{24pt}\selectfont \paperTitle\par}
\end{center}

\begin{center}
    {\FZhei\zihao{4}\setlength{\baselineskip}{24pt}\selectfont 摘要\par}
\end{center}

本文面向人口老龄化背景下社区养老服务配置不均衡的问题，以福州市主城区为研究对象，旨在识别社区养老服务供给与老年需求之间的空间错配，并提出有限资源约束下的服务补位方案。研究基于人口栅格、养老服务设施、医疗设施、代际共享服务设施、道路交通与公共服务 POI 等多源数据，构建 500\,m $\times$ 500\,m 网格化分析单元；在方法上，首先估计老年服务需求强度，其次构建代际共享修正的 2SFCA 模型测算养老服务可达性，再利用 GraphSAGE 识别空间关联下的高错配风险网格，并结合 LightGBM-SHAP 解释错配成因，最后采用 MCLP 最大覆盖模型进行服务补位优化。结果部分应根据实证结果填写高错配区域、主要影响因素和补位效果。研究创新在于将代际共享服务纳入社区养老可达性修正框架，并把错配识别、机器学习解释和设施布局优化连接起来，为福州市代际友好型社区建设和社区养老服务精准配置提供量化依据。

\keywords{社区养老；供需错配；代际友好；可达性；福州市}

\clearpage

% ================================================================
% 章节目录：引用至正文第三级标题
% ================================================================
\phantomsection
\tableofcontents
\clearpage

% ================================================================
% 表格与插图清单
% ================================================================
\listoftablesandfigures
\clearpage

% ================================================================
% 正文：正文第一页开始编页码
% ================================================================
\pagenumbering{arabic}
\setcounter{page}{1}
\pagestyle{plain}
\maintextformat

\begin{center}
    {\xiaobiaosong\zihao{3}\setlength{\baselineskip}{24pt}\selectfont \paperTitle\par}
\end{center}

\section{引言}

\subsection{研究背景}

\subsubsection{人口老龄化与社区养老服务压力}

说明我国老龄化加深，养老服务需求从机构养老逐渐转向社区居家养老、助餐、医养结合和日常生活支持。

这一节不要写太长，主要作用是引出问题。

\subsubsection{福州市社区养老服务的现实基础}

说明福州既有老城区老龄人口集聚，也有新区扩张和城乡养老资源差异。福州适合作为研究对象，因为既有养老服务设施，也有较强的智慧养老和社区服务建设基础。

\subsubsection{从“设施覆盖”到“有效可达”的研究转向}

强调本文不是简单统计养老设施数量，而是研究老人是否能真正获得服务。这里可以提出核心判断：

\begin{quote}
养老服务设施存在，不等于养老服务可达；设施覆盖率高，也不等于供需匹配。
\end{quote}

\subsection{研究问题}

\subsubsection{福州市老年服务需求如何空间分布}

回答“老人在哪里”。

\subsubsection{社区养老服务是否与老年需求匹配}

回答“服务是否供给到真正需要的地方”。

\subsubsection{代际共享服务能否缓解养老服务压力}

说明社区食堂、公园、图书馆、社区活动中心、体育空间、便民服务中心等设施虽然不是传统养老设施，但可以为老年群体提供助餐、休闲、社交和生活支持。你们前期材料已经把社区服务划分为银发服务、青年活力和代际共享服务三类，这里可以作为指标体系的来源。

\subsubsection{有限资源下应优先补什么、补在哪里}

回答“补位优化”的问题。

\subsection{研究思路与技术路线}

\subsubsection{总体研究思路}

可以写成：

\begin{quote}
本文按照“需求识别—服务可达—错配诊断—成因解释—智能补位”的逻辑展开。
\end{quote}

\subsubsection{技术路线设计}

建议放一张技术路线图。

图中包括：

数据采集 → 网格化处理 → 老年需求估计 → 代际共享修正 2SFCA → 供需错配识别 → GraphSAGE 风险识别 → LightGBM-SHAP 成因解释 → MCLP 服务补位优化。

\subsubsection{研究创新}

写三点即可。

第一，构建代际共享修正的养老服务可达性测度方法。
第二，引入图神经网络识别空间关联下的高错配风险网格。
第三，将错配识别与 MCLP 补位优化结合，形成可执行的设施布局建议。

\section{数据来源与变量构建}

\subsection{研究区域与空间单元}

\subsubsection{研究区域选择}

研究对象建议设定为福州市主城区，包括鼓楼区、台江区、仓山区、晋安区、马尾区。如数据允许，可扩展至长乐区和闽侯县核心城镇片区。

\subsubsection{空间网格划分}

将研究区域划分为 500\,m $\times$ 500\,m 网格。每个网格作为基本分析单元。

说明为什么不用行政街道：

\begin{quote}
街道尺度过粗，难以刻画社区服务半径；网格尺度更适合计算服务可达性和补位点位。
\end{quote}

\subsubsection{数据空间匹配}

将人口、养老设施、医疗设施、共享服务设施、交通设施全部匹配到网格中。

这里可以写：

\begin{quote}
对所有 POI 数据进行坐标统一、去重、分类和网格归属处理。
\end{quote}

\subsection{数据来源}

\subsubsection{人口与社会经济数据}

包括：福州市统计年鉴、人口普查数据、区县人口结构、老龄化比例、常住人口、GDP、居民收入、财政支出等。

用途：估计老年需求强度，作为模型控制变量。

\subsubsection{养老服务设施数据}

包括：养老机构、护理院、老年公寓、长者食堂、日间照料中心、养老服务站、农村幸福院、助餐点等。

来源：福州养老服务资源地图、高德 POI、百度 POI、民政部门公开数据。

用途：构建养老服务供给能力。

\subsubsection{医养协同服务数据}

包括：医院、社区卫生服务中心、药店、康复机构、长护险定点机构。

用途：衡量老年群体获得医疗和护理支持的能力。

\subsubsection{代际共享服务数据}

包括：社区食堂、平价餐饮、公园、图书馆、体育空间、社区活动中心、便民服务中心、公交站、地铁站。

用途：衡量社区是否具有可同时服务老人和其他年龄群体的共享型服务环境。

\subsubsection{空间交通数据}

包括：道路网络、公交站、地铁站、行政边界、步行距离、公交出行时间等。

用途：计算养老服务可达性和服务圈覆盖。

\subsection{指标体系设计}

\subsubsection{老年需求指标}

\begin{table}[H]
\centering
\caption{老年需求指标体系}
\begin{tabular}{ll}
\toprule
\rowcolor{tablehead}
指标 & 含义 \\
\midrule
老年人口密度 & 基础养老需求 \\
高龄人口密度 & 护理型养老需求 \\
人口密度 & 社区服务总体压力 \\
住宅小区密度 & 居住人口集聚程度 \\
老旧小区密度 & 潜在老龄人口集聚 \\
到医院距离 & 健康服务压力 \\
\bottomrule
\end{tabular}
\end{table}

说明：这一组指标回答“哪里更需要养老服务”。

\subsubsection{养老服务供给指标}

\begin{table}[H]
\centering
\caption{养老服务供给指标体系}
\begin{tabular}{ll}
\toprule
\rowcolor{tablehead}
指标 & 含义 \\
\midrule
养老机构数量 & 机构养老供给 \\
长者食堂数量 & 助餐服务供给 \\
日间照料中心数量 & 社区照护服务 \\
养老服务站数量 & 居家养老支持 \\
护理机构数量 & 失能照护供给 \\
养老床位数 & 机构容量，若缺失可用机构类型权重代替 \\
\bottomrule
\end{tabular}
\end{table}

说明：这一组指标回答“养老服务在哪里”。

\subsubsection{医养协同指标}

\begin{table}[H]
\centering
\caption{医养协同指标体系}
\begin{tabular}{ll}
\toprule
\rowcolor{tablehead}
指标 & 含义 \\
\midrule
医院密度 & 高等级医疗支持 \\
社区卫生服务中心密度 & 基层医疗支持 \\
药店密度 & 日常健康服务 \\
康复机构密度 & 康复照护支持 \\
长护险定点机构密度 & 护理型服务能力 \\
\bottomrule
\end{tabular}
\end{table}

说明：这一组指标回答“养老服务是否具备医疗支撑”。

\subsubsection{代际共享服务指标}

\begin{table}[H]
\centering
\caption{代际共享服务指标体系}
\begin{tabular}{ll}
\toprule
\rowcolor{tablehead}
指标 & 含义 \\
\midrule
社区食堂 / 平价餐饮 & 助餐与日常消费 \\
公园 / 广场 & 休闲活动 \\
图书馆 / 文化馆 & 公共文化服务 \\
体育空间 & 健康活动 \\
社区活动中心 & 社交互动 \\
便民服务中心 & 生活便利 \\
公交 / 地铁站 & 交通可达 \\
\bottomrule
\end{tabular}
\end{table}

说明：这一组指标回答“哪些服务能为养老压力提供间接缓冲”。

\section{模型构建}

\subsection{老年需求空间估计模型}

\subsubsection{人口栅格下推思路}

由于网格尺度老年人口数据通常无法直接获得，本文采用区县老龄化比例、人口栅格和居住 POI 进行下推估计。

\subsubsection{老年需求强度计算}

设网格 $g$ 属于区县 $c$，其老年需求强度为：

\begin{equation}
D_g = \operatorname{Pop}_g \times r_c^{65+} \times W_g
\end{equation}

其中，$\operatorname{Pop}_g$ 是网格人口，$r_c^{65+}$ 是区县 65 岁及以上人口比例，$W_g$ 是居住修正权重。

\subsubsection{居住修正权重设定}

居住修正权重由住宅小区密度、老旧小区密度和建成区强度构成。

说明：这一节不需要展开过深，重点是说明需求估计不是简单平均分配，而是考虑居住空间差异。

\subsection{养老服务有效供给测度}

\subsubsection{服务设施分类赋权}

将养老服务设施划分为机构养老、社区照护、助餐服务、医养协同和代际共享五类。

建议设定：

\begin{table}[H]
\centering
\caption{养老服务类型与权重含义}
\begin{tabular}{ll}
\toprule
\rowcolor{tablehead}
服务类型 & 权重含义 \\
\midrule
养老机构 & 机构照护能力强 \\
日间照料中心 & 社区照护能力较强 \\
养老服务站 & 居家服务支持 \\
长者食堂 & 助餐支持 \\
医疗与药店 & 医养协同支持 \\
代际共享设施 & 生活与社交支持 \\
\bottomrule
\end{tabular}
\end{table}

\subsubsection{服务供给指数计算}

构建网格服务供给指数：

\begin{equation}
S_g = \sum_{k=1}^{m} w_k X_{gk}
\end{equation}

其中，$X_{gk}$ 是网格 $g$ 中第 $k$ 类服务指标，$w_k$ 是对应权重。

\subsubsection{指标标准化处理}

所有指标统一进行极差标准化或 Z-score 标准化，避免不同量纲影响结果。

说明：这里不要把 CRITIC-TOPSIS 写太重。供给指数只是为 2SFCA 提供输入，主角不是综合评价。

\subsection{代际共享修正 2SFCA 可达性模型}

\subsubsection{服务点供需比计算}

对每个服务点 $j$，计算其服务范围内的供需比：

\begin{equation}
R_j = \frac{S_j}{\sum_{g \in d_{gj} \leq d_0} D_g f(d_{gj})}
\end{equation}

其中，$S_j$ 是服务点有效供给能力，$D_g$ 是网格需求，$d_{gj}$ 是网格到服务点的出行时间。

\subsubsection{网格服务可达性计算}

对每个网格 $g$，计算其可获得的养老服务水平：

\begin{equation}
A_g = \sum_{j \in d_{gj} \leq d_0} R_j f(d_{gj})
\end{equation}

其中，$A_g$ 越大，说明养老服务可达性越好。

\subsubsection{距离衰减函数设定}

采用高斯距离衰减：

\begin{equation}
f(d)=e^{-(d/\beta)^2}
\end{equation}

说明：距离越远，服务可获得性越低。这比简单 15 分钟圈更合理。

\subsubsection{代际共享服务修正}

构建共享服务支撑项：

\begin{equation}
A_g^{*}=A_g+\alpha \operatorname{Shared}_g
\end{equation}

其中，$\operatorname{Shared}_g$ 表示网格代际共享服务水平，$\alpha$ 表示共享服务对养老可达性的缓冲系数。

建议正文设：

\begin{equation}
\alpha = 0.2
\end{equation}

在稳健性中检验：

\begin{equation}
\alpha = 0.1,\ 0.2,\ 0.3
\end{equation}

说明：这是你们论文的第一个亮点。它体现了社区养老不是孤立的养老设施问题，而是与社区共享服务体系有关。

\subsection{供需错配指数构建}

\subsubsection{基础错配指数}

构建：

\begin{equation}
\operatorname{Mismatch}_g = Z(D_g) - Z(A_g^{*})
\end{equation}

其中，$Z(D_g)$ 是标准化需求，$Z(A_g^{*})$ 是修正后可达性。

\subsubsection{错配程度分级}

将 $\operatorname{Mismatch}_g$ 按分位数划分为：

\begin{table}[H]
\centering
\caption{错配程度分级标准}
\begin{tabular}{ll}
\toprule
\rowcolor{tablehead}
分位 & 等级 \\
\midrule
前 20\% & 高错配 \\
20\%—80\% & 中错配 \\
后 20\% & 低错配 \\
\bottomrule
\end{tabular}
\end{table}

\subsubsection{供需错配四象限}

根据需求强度和可达性将网格划分为：

\begin{table}[H]
\centering
\caption{供需错配四象限类型}
\begin{tabular}{ll}
\toprule
\rowcolor{tablehead}
类型 & 含义 \\
\midrule
高需求—高供给 & 基本匹配区 \\
高需求—低供给 & 优先补位区 \\
低需求—高供给 & 可能冗余区 \\
低需求—低供给 & 基础薄弱区 \\
\bottomrule
\end{tabular}
\end{table}

说明：这一节直接服务后面的政策建议。

\subsection{GraphSAGE 空间错配识别模型}

\subsubsection{空间图构建}

将 500\,m 网格视为节点，若两个网格空间相邻，则建立边。

节点特征包括：

\begin{itemize}
\item 老年需求强度；
\item 养老服务可达性；
\item 医养协同水平；
\item 代际共享服务水平；
\item 交通可达性；
\item 社区活力指标。
\end{itemize}

\subsubsection{高错配风险标签构造}

根据 $\operatorname{Mismatch}_g$ 的分位数将网格划分为高、中、低错配风险。

GraphSAGE 的任务是预测网格错配风险等级。

\subsubsection{模型聚合思想}

GraphSAGE 通过聚合邻近网格的特征来更新节点表示：

\begin{equation}
h_v^k = \sigma \left(W^k \cdot \operatorname{CONCAT}\left(h_v^{k-1}, \operatorname{AGG}\{h_u^{k-1}\mid u \in N(v)\}\right)\right)
\end{equation}

说明：正文中不需要讲太多神经网络原理。重点写它的作用：

\begin{quote}
GraphSAGE 用于识别考虑邻近社区资源共享后的高错配风险区域。
\end{quote}

\subsubsection{模型输出}

输出：

\begin{itemize}
\item 高错配风险概率；
\item 高风险网格空间分布；
\item 与传统错配指数的对比结果。
\end{itemize}

\subsection{LightGBM-SHAP 成因解释模型}

\subsubsection{被解释变量}

被解释变量为：

\begin{equation}
\operatorname{Mismatch}_g
\end{equation}

或高错配风险等级。

\subsubsection{解释变量}

包括：

\begin{itemize}
\item 老年人口密度；
\item 高龄人口密度；
\item 长者食堂密度；
\item 养老服务站密度；
\item 医疗设施密度；
\item 公共交通可达性；
\item 代际共享服务密度；
\item 道路密度；
\item 社区活力指标。
\end{itemize}

\subsubsection{LightGBM 建模}

使用 LightGBM 拟合非线性关系，识别影响错配的关键变量。

\subsubsection{SHAP 解释}

利用 SHAP 输出变量贡献。

重点分析：

\begin{itemize}
\item 哪些变量推高错配；
\item 哪些变量缓解错配；
\item 不同区域的主要短板是否不同。
\end{itemize}

说明：这一节主要靠 SHAP 图说话，文字不宜过长。

\subsection{MCLP 服务补位优化模型}

\subsubsection{候选补位点设定}

候选点包括：

\begin{itemize}
\item 社区服务中心；
\item 居委会；
\item 社区卫生服务中心；
\item 现有养老服务站；
\item 长者食堂周边；
\item 高错配网格中心。
\end{itemize}

\subsubsection{最大覆盖目标函数}

目标是在有限新增点位下覆盖最多高需求老人：

\begin{equation}
\max \sum_g D_g y_g
\end{equation}

约束：

\begin{equation}
y_g \leq \sum_{j \in N_g}x_j
\end{equation}

\begin{equation}
\sum_j x_j \leq p
\end{equation}

其中，$x_j=1$ 表示在候选点 $j$ 新建或升级服务点，$y_g=1$ 表示网格 $g$ 被覆盖。

\subsubsection{政策情景设定}

设置三种情景：

\begin{table}[H]
\centering
\caption{MCLP 政策情景设定}
\begin{tabular}{ll}
\toprule
\rowcolor{tablehead}
情景 & 新增 / 升级点位 \\
\midrule
保守情景 & 10 个 \\
中等情景 & 20 个 \\
积极情景 & 30 个 \\
\bottomrule
\end{tabular}
\end{table}

\subsubsection{优化效果评价}

评价指标包括：

\begin{itemize}
\item 高需求老人覆盖率；
\item 平均可达时间；
\item 高错配网格数量；
\item 服务公平性变化。
\end{itemize}

说明：MCLP 是你们论文的落地模型。它负责回答“补哪里”。

\section{实证结果分析}

\subsection{福州市老年需求空间分布}

\subsubsection{老年需求强度整体格局}

展示老年需求地图。说明哪些区域老年需求较高。

\subsubsection{高龄照护需求分布}

展示高龄需求或护理需求分布。重点分析是否存在高龄老人集聚片区。

\subsubsection{老年需求热点区识别}

列出需求强度排名前若干网格或街道。

\subsection{社区养老服务供给格局}

\subsubsection{养老服务设施空间分布}

展示养老机构、长者食堂、日间照料中心、养老服务站的整体分布。

\subsubsection{医养协同服务分布}

展示医院、社区卫生服务中心、药店、康复机构等分布。

\subsubsection{代际共享服务分布}

展示公园、社区食堂、公共交通、文化体育空间等分布。

\subsection{养老服务可达性测度结果}

\subsubsection{15 分钟社区养老服务可达性}

展示长者食堂、养老服务站等近距离服务可达性。

\subsubsection{30 分钟照护服务可达性}

展示日间照料中心、社区照护设施可达性。

\subsubsection{代际共享修正前后对比}

展示修正前后的可达性差异。

说明：这一节要突出你们的创新：共享服务确实会改变部分区域的服务评价。

\subsection{供需错配识别结果}

\subsubsection{错配指数空间分布}

展示 $\operatorname{Mismatch}_g$ 地图。说明高错配区集中在哪里。

\subsubsection{四象限分类结果}

展示高需求—低供给、高需求—高供给等四类区域。

\subsubsection{优先补位区识别}

列出高需求—低供给区域，并说明其主要短板。

\subsection{GraphSAGE 高错配风险识别结果}

\subsubsection{模型识别结果}

展示 GraphSAGE 输出的高风险网格图。

\subsubsection{与传统错配指数对比}

比较 GraphSAGE 与基础错配指数结果是否一致。

重点说明：GraphSAGE 能够识别部分受邻近区域影响的潜在风险网格。

\subsubsection{高风险区域特征分析}

总结高风险区域通常具有哪些特征，例如：

\begin{itemize}
\item 老年需求高；
\item 养老服务弱；
\item 医疗协同不足；
\item 公共交通不便；
\item 共享服务支撑不足。
\end{itemize}

\subsection{LightGBM-SHAP 成因解释结果}

\subsubsection{模型拟合效果}

简要说明模型精度。不要占太多篇幅。

\subsubsection{全局特征重要性}

展示 SHAP 特征重要性图。说明哪些变量最影响错配。

\subsubsection{不同短板类型的成因差异}

分析不同区域的主要短板：

\begin{itemize}
\item 助餐不足型；
\item 医养协同不足型；
\item 交通可达性不足型；
\item 共享服务不足型；
\item 综合服务不足型。
\end{itemize}

\subsection{MCLP 服务补位优化结果}

\subsubsection{新增 10 个服务点情景}

展示保守情景下补位点位图。说明覆盖了哪些高需求区域。

\subsubsection{新增 20 个服务点情景}

展示中等情景下补位点位图。说明覆盖效果提升。

\subsubsection{新增 30 个服务点情景}

展示积极情景下补位点位图。说明边际收益是否下降。

\subsubsection{补位前后效果对比}

放一张表：

\begin{table}[H]
\centering
\caption{补位前后效果对比指标}
\begin{tabular}{lllll}
\toprule
\rowcolor{tablehead}
指标 & 补位前 & 新增 10 个 & 新增 20 个 & 新增 30 个 \\
\midrule
覆盖老年人口 &  &  &  &  \\
平均可达时间 &  &  &  &  \\
高错配网格数量 &  &  &  &  \\
服务公平性指数 &  &  &  &  \\
\bottomrule
\end{tabular}
\end{table}

\section{稳健性检验与讨论}

\subsection{服务半径敏感性检验}

\subsubsection{15 分钟阈值调整}

将长者食堂和社区服务站阈值调整为 10 分钟、20 分钟。

\subsubsection{30 分钟阈值调整}

将日间照料服务阈值调整为 20 分钟、40 分钟。

\subsubsection{结果稳定性分析}

说明高错配区域是否保持稳定。

\subsection{代际共享修正系数敏感性检验}

\subsubsection{不同共享修正系数取值设定}

检验：

\begin{equation}
\alpha = 0.1,\ 0.2,\ 0.3
\end{equation}

\subsubsection{高错配区域变化}

比较不同 $\alpha$ 下的高错配网格是否变化明显。

\subsubsection{稳健性结论}

若主要高错配区域基本一致，说明模型稳定。

\subsection{模型对比检验}

\subsubsection{LightGBM 与传统模型对比}

可与随机森林、逻辑回归或 XGBoost 简单对比。

\subsubsection{GraphSAGE 与非图模型对比}

说明加入空间邻接信息后，识别效果是否提升。

\subsubsection{检验结论}

这一节不需要长篇，只需说明主结论稳健。

\section{政策建议}

\subsection{高龄照护缺口区：补护理型服务}

\subsubsection{区域特征}

高龄人口密度高，医疗和护理服务不足。

\subsubsection{补位方向}

新增护理型社区养老服务点、长护险定点服务、社区康复和上门护理。

\subsubsection{实施建议}

优先依托社区卫生服务中心和现有养老服务站升级。

\subsection{助餐服务不足区：补长者食堂和社区食堂}

\subsubsection{区域特征}

老年需求高，但长者食堂和助餐点不足。

\subsubsection{补位方向}

新增长者食堂、社区食堂、移动助餐和配送服务。

\subsubsection{实施建议}

优先布局在老旧小区和高需求网格周边。

\subsection{居家照料不足区：补日间照料中心和养老服务站}

\subsubsection{区域特征}

机构养老不一定缺，但社区居家服务弱。

\subsubsection{补位方向}

新增日间照料中心、社区养老服务站、助浴助洁助医服务。

\subsubsection{实施建议}

推动服务站向综合照护站升级。

\subsection{代际共享服务短板区：补复合型公共服务空间}

\subsubsection{区域特征}

养老服务需求高，同时缺少公园、体育空间、图书馆、社区活动中心等共享服务。

\subsubsection{补位方向}

补充社区食堂、公共活动空间、文体设施和便民服务中心。

\subsubsection{实施建议}

推动“一站多用”的复合型社区服务空间建设。

\subsection{供给冗余区：推动存量设施转型}

\subsubsection{区域特征}

服务供给较多，但老年需求相对不足。

\subsubsection{优化方向}

控制新增设施，提升服务质量，推动存量设施向助餐、康复、文体和上门服务转型。

\subsubsection{实施建议}

避免重复建设，提高现有设施利用率。

\section{结论与不足}

\subsection{主要研究结论}

\subsubsection{福州市社区养老服务存在空间差异}

总结老年需求和养老服务供给的空间分布差异。

\subsubsection{高需求低供给区域具有明显识别价值}

总结哪些类型区域需要优先补位。

\subsubsection{代际共享服务能够部分缓解养老服务压力}

总结共享服务修正后的变化。

\subsubsection{MCLP 能为服务点布局提供量化依据}

总结新增 10、20、30 个点位的优化效果。

\subsection{研究贡献}

\subsubsection{方法贡献}

提出代际共享修正的养老服务可达性测度方法。

\subsubsection{实践贡献}

为福州市社区养老服务补位提供可执行方案。

\subsubsection{政策贡献}

为代际友好型社区建设提供量化依据。

\subsection{研究不足}

\subsubsection{老年人口网格估计存在误差}

由于缺少社区级真实老年人口，只能通过人口栅格和区县比例估计。

\subsubsection{部分服务设施容量数据可能缺失}

如床位数、收费标准、服务人数等字段不完整。

\subsubsection{模型结果仍需结合实地调研验证}

补位点位应结合土地、财政、社区意愿和实际运营条件进一步评估。

% ================================================================
% 参考文献：按先中文、后英文顺序列出，正式提交前请替换为实际引用文献
% ================================================================
\clearpage
\phantomsection
\addcontentsline{toc}{section}{参考文献}

\begin{center}
    {\FZhei\zihao{4}\setlength{\baselineskip}{24pt}\selectfont 参考文献\par}
\end{center}

\begin{enumerate}[label={\textnormal{[\arabic*]}},leftmargin=2em]
    \item 作者. 中文文献题名[J]. 期刊名, 年份, 卷(期): 起止页码.
    \item 作者. 中文专著题名[M]. 出版地: 出版社, 年份.
    \item Author A, Author B. English article title[J]. Journal Name, Year, Volume(Issue): Page--Page.
\end{enumerate}

% ================================================================
% 附录：正文主体的补充说明
% ================================================================
\clearpage
\phantomsection
\addcontentsline{toc}{section}{附录}

\begin{center}
    {\FZhei\zihao{4}\setlength{\baselineskip}{24pt}\selectfont 附录\par}
\end{center}

{\FZkai\zihao{4}\setlength{\baselineskip}{24pt}\selectfont 附录1\quad POI 分类规则表\par}
\addcontentsline{toc}{subsection}{附录1\quad POI 分类规则表}

列出养老服务、医养协同、代际共享服务、社区活力 POI 的关键词、纳入标准、剔除规则和分类方式。

{\FZkai\zihao{4}\setlength{\baselineskip}{24pt}\selectfont 附录2\quad 变量标准化方法\par}
\addcontentsline{toc}{subsection}{附录2\quad 变量标准化方法}

说明极差标准化、Z-score 标准化以及异常值处理方法。

{\FZkai\zihao{4}\setlength{\baselineskip}{24pt}\selectfont 附录3\quad GraphSAGE 模型参数\par}
\addcontentsline{toc}{subsection}{附录3\quad GraphSAGE 模型参数}

列出隐藏层维度、学习率、训练轮数、邻居采样数、训练集与测试集划分方法等。

{\FZkai\zihao{4}\setlength{\baselineskip}{24pt}\selectfont 附录4\quad MCLP 候选点清单\par}
\addcontentsline{toc}{subsection}{附录4\quad MCLP 候选点清单}

列出候选点来源、筛选规则、服务半径设定和不同情景下的候选点选择结果。

% ================================================================
% 致谢：内容应简洁明了、实事求是
% ================================================================
\clearpage
\phantomsection
\addcontentsline{toc}{section}{致谢}

\begin{center}
    {\FZhei\zihao{4}\setlength{\baselineskip}{24pt}\selectfont 致谢\par}
\end{center}

论文的完成离不开指导教师、相关课程教师以及团队成员在选题、数据收集、模型构建和论文修改过程中的帮助。感谢所有为本研究提供建议、资料和支持的老师、同学及有关单位。本文仍有不足之处，恳请各位专家批评指正。

\vspace{2em}

\begin{flushright}
2026年\quad 月\quad 日
\end{flushright}

\end{document}